%==============================================================================
\section{Thuật toán K-Means}
\label{sec:k-means}
%==============================================================================

\begin{dinhnghia}[title={K-Means Clustering}]
Thuật toán K-Means tính centroid và lặp cho đến khi đạt centroid tối ưu. Giả định \textbf{đã biết trước} số cụm; còn gọi là thuật toán phân cụm \textbf{phẳng} (flat clustering). Ký hiệu $K$ là số cụm mà thuật toán xác định từ dữ liệu.
\end{dinhnghia}

\begin{ynghia}[title={Mục tiêu tối thiểu}]
Mỗi điểm được gán vào cụm sao cho \textbf{tổng bình phương khoảng cách} giữa điểm và centroid là nhỏ nhất. Biến thiên trong cụm càng nhỏ thì các điểm trong cùng cụm càng giống nhau.
\end{ynghia}

\subsection{Cách hoạt động của K-Means}

\begin{phuongphap}[title={Các bước}]
\begin{itemize}
  \item \textbf{Bước 1}: Chỉ định số cụm $K$ cần sinh.
  \item \textbf{Bước 2}: Chọn ngẫu nhiên $K$ điểm và gán ban đầu (phân loại theo số điểm khởi tạo).
  \item \textbf{Bước 3}: Tính centroid của từng cụm.
  \item \textbf{Bước 4}: Lặp đến khi phân gán ổn định:
  \begin{itemize}
    \item \textbf{4.1}: Tính tổng bình phương khoảng cách giữa điểm và centroid.
    \item \textbf{4.2}: Gán mỗi điểm vào cụm có centroid \textbf{gần hơn} các cụm còn lại.
    \item \textbf{4.3}: Cập nhật centroid bằng \textbf{trung bình} mọi điểm trong cụm.
  \end{itemize}
\end{itemize}
\end{phuongphap}

\begin{dinhnghia}[title={Expectation--Maximization (EM)}]
K-means theo khung EM:
\begin{itemize}
  \item \textbf{Bước E (Expectation - Kỳ vọng)}: gán điểm vào cụm gần nhất.
  \item \textbf{Bước M (Maximization - Cực đại hóa)}: tính lại centroid mỗi cụm.
\end{itemize}
\end{dinhnghia}

\begin{luuy}[title={Chuẩn hóa dữ liệu}]
Thuật toán dùng đo \textbf{khoảng cách} để đo tương tự; nên \textbf{chuẩn hóa} (standardize) feature trước khi chạy K-means và các thuật toán phân cụm tương tự.
\end{luuy}

\begin{luuy}[title={Cục cực địa phương}]
Do khởi tạo centroid ngẫu nhiên và lặp lặp lại, K-means có thể kẹt ở \textbf{cục cực địa phương} thay vì toàn cục. Nên thử \textbf{nhiều khởi tạo} centroid khác nhau (ví dụ \texttt{n\_init} trong scikit-learn).
\end{luuy}

\begin{tinhchat}[title={Ưu/nhược điểm tổng quan}]
Thuật toán đơn giản, hiệu quả với dataset lớn; nhưng nhạy khởi tạo, dễ hội tụ cục bộ, và thường giả định phương sai đồng đều giữa các cụm.
\end{tinhchat}

\subsection{Mục tiêu phân tích cụm}

\begin{dongco}[title={Hai mục tiêu chính}]
\begin{itemize}
  \item Rút ra \textbf{trực giác có ý nghĩa} từ dữ liệu.
  \item \textbf{Phân cụm rồi dự đoán}: xây mô hình riêng cho từng nhóm con sau khi phân cụm.
\end{itemize}
\end{dongco}

\subsection{Triển khai K-Means bằng Python}

\begin{phuongphap}[title={scikit-learn}]
Python có scikit-learn với lớp \texttt{KMeans} trong \texttt{sklearn.cluster}.
\end{phuongphap}

\subsubsection{Ví dụ 1 --- Hiểu K-Means với dữ liệu ngẫu nhiên}

\begin{dongco}[title={Bối cảnh}]
Sinh 300 điểm hai chiều ngẫu nhiên, áp dụng K-means với $K=3$.
\end{dongco}

\begin{vidu}[title={Bước 1 --- Import}]
\begin{lstlisting}
import numpy as np
import matplotlib.pyplot as plt
from sklearn.cluster import KMeans
\end{lstlisting}
\end{vidu}

\begin{vidu}[title={Bước 2 --- Sinh và vẽ dữ liệu}]
\begin{lstlisting}
X = np.random.rand(300, 2)

plt.figure(figsize=(7.5, 3.5))
plt.scatter(X[:, 0], X[:, 1], s=20)
plt.show()
\end{lstlisting}
\end{vidu}

\begin{ynghia}[title={Dữ liệu ban đầu}]
\tsimg{03_k_means/img01.jpg}
\end{ynghia}

\begin{vidu}[title={Bước 3 --- Khởi tạo K-Means}]
\begin{lstlisting}
kmeans = KMeans(n_clusters=3, max_iter=100)
\end{lstlisting}
\end{vidu}

\begin{vidu}[title={Bước 4 --- Huấn luyện}]
\begin{lstlisting}
kmeans.fit(X)
\end{lstlisting}
\end{vidu}

\begin{vidu}[title={Bước 5 --- Vẽ cụm và centroid}]
\begin{lstlisting}
plt.figure(figsize=(7.5, 3.5))
plt.scatter(X[:, 0], X[:, 1], c=kmeans.labels_, s=20)
plt.scatter(kmeans.cluster_centers_[:, 0],
            kmeans.cluster_centers_[:, 1],
            marker='x', c='r', s=50, alpha=0.9)
plt.show()
\end{lstlisting}
\end{vidu}

\begin{ynghia}[title={Kết quả phân cụm}]
Điểm tô màu theo nhãn cụm; centroid đánh dấu \texttt{x} màu đỏ.

\tsimg{03_k_means/img02.jpg}
\end{ynghia}

\begin{vidu}[title={Code đầy đủ (ví dụ 1)}]
\begin{lstlisting}
import numpy as np
import matplotlib.pyplot as plt
from sklearn.cluster import KMeans

X = np.random.rand(300, 2)

kmeans = KMeans(n_clusters=3, max_iter=100)
kmeans.fit(X)

plt.figure(figsize=(7.5, 3.5))
plt.scatter(X[:, 0], X[:, 1], c=kmeans.labels_, s=20)
plt.scatter(kmeans.cluster_centers_[:, 0],
            kmeans.cluster_centers_[:, 1],
            marker='x', c='r', s=50, alpha=0.9)
plt.show()
\end{lstlisting}
\end{vidu}

\subsubsection{Ví dụ 2 --- Dataset 2D bốn blob}

\begin{dongco}[title={Bối cảnh}]
Sinh dataset 2D gồm 4 blob bằng \texttt{make\_blobs}, huấn luyện K-means với 4 cụm.
\end{dongco}

\begin{vidu}[title={Import và sinh blob}]
\begin{lstlisting}
import matplotlib.pyplot as plt
import numpy as np
from sklearn.cluster import KMeans
from sklearn.datasets import make_blobs

X, y_true = make_blobs(n_samples=400, centers=4,
                       cluster_std=0.60, random_state=0)
\end{lstlisting}
\end{vidu}

\begin{vidu}[title={Vẽ dữ liệu gốc}]
\begin{lstlisting}
plt.scatter(X[:, 0], X[:, 1], s=20)
plt.show()
\end{lstlisting}
\end{vidu}

\begin{ynghia}[title={Trực quan 2D blob}]
\tsimg{03_k_means/img03.jpg}
\end{ynghia}

\begin{vidu}[title={Huấn luyện và dự đoán}]
\begin{lstlisting}
kmeans = KMeans(n_clusters=4)
kmeans.fit(X)
y_kmeans = kmeans.predict(X)
\end{lstlisting}
\end{vidu}

\begin{vidu}[title={Vẽ cụm và tâm}]
\begin{lstlisting}
plt.scatter(X[:, 0], X[:, 1], c=y_kmeans, s=20)
centers = kmeans.cluster_centers_
plt.scatter(centers[:, 0], centers[:, 1],
            c='blue', s=100, alpha=0.9)
plt.show()
\end{lstlisting}
\end{vidu}

\begin{ynghia}[title={Centroid học được}]
\tsimg{03_k_means/img04.jpg}
\end{ynghia}

\subsubsection{Ví dụ 3 --- Chữ số viết tay (digits)}

\begin{dongco}[title={Bối cảnh}]
Áp dụng K-means lên \texttt{load\_digits}: tìm nhóm chữ số tương tự \textbf{không} dùng nhãn gốc khi gom cụm.
\end{dongco}

\begin{vidu}[title={Nạp dữ liệu}]
\begin{lstlisting}
from sklearn.datasets import load_digits

digits = load_digits()
digits.data.shape   # (1797, 64)
\end{lstlisting}
\end{vidu}

\begin{ynghia}[title={Kích thước}]
1797 mẫu, mỗi mẫu 64 feature (ảnh $8 \times 8$).
\end{ynghia}

\begin{vidu}[title={Phân cụm 10 nhóm}]
\begin{lstlisting}
kmeans = KMeans(n_clusters=10, random_state=0)
clusters = kmeans.fit_predict(digits.data)
kmeans.cluster_centers_.shape   # (10, 64)
\end{lstlisting}
\end{vidu}

\begin{vidu}[title={Vẽ tâm cụm dạng ảnh}]
\begin{lstlisting}
fig, ax = plt.subplots(2, 5, figsize=(8, 3))
centers = kmeans.cluster_centers_.reshape(10, 8, 8)
for axi, center in zip(ax.flat, centers):
    axi.set(xticks=[], yticks=[])
    axi.imshow(center, interpolation='nearest',
               cmap=plt.cm.binary)
\end{lstlisting}
\end{vidu}

\begin{ynghia}[title={Centroid mỗi cụm}]
\tsimg{03_k_means/img05.jpg}
\end{ynghia}

\begin{phuongphap}[title={Ghép nhãn cụm với nhãn thật}]
Dùng \texttt{scipy.stats.mode} để gán mỗi cụm nhãn digit phổ biến nhất trong cụm đó, rồi tính \texttt{accuracy\_score}.
\end{phuongphap}

\begin{vidu}[title={Gán nhãn và độ chính xác}]
\begin{lstlisting}
from scipy.stats import mode
from sklearn.metrics import accuracy_score

labels = np.zeros_like(clusters)
for i in range(10):
    mask = (clusters == i)
    labels[mask] = mode(digits.target[mask], keepdims=True).mode[0]

accuracy_score(digits.target, labels)   # ~0.79
\end{lstlisting}
\end{vidu}

\begin{tinhchat}[title={Độ chính xác ~80\%}]
Ghép nhãn cụm với nhãn gốc cho độ chính xác khoảng \textbf{80\%} --- minh họa cụm bắt được cấu trúc chữ số dù không dùng nhãn lúc fit.
\end{tinhchat}

\subsection{Ưu điểm K-Means}

\begin{tinhchat}[title={Danh sách ưu điểm}]
\begin{itemize}
  \item Dễ hiểu và triển khai.
  \item Với số biến lớn, thường \textbf{nhanh hơn} phân cụm phân cấp.
  \item Khi tính lại centroid, một mẫu có thể \textbf{đổi cụm}.
  \item Cụm thường \textbf{chặt} hơn so với phân cụm phân cấp trong nhiều bài toán.
\end{itemize}
\end{tinhchat}

\subsection{Nhược điểm K-Means}

\begin{luuy}[title={Hạn chế}]
\begin{itemize}
  \item Khó chọn $k$ (số cụm).
  \item Kết quả phụ thuộc mạnh vào $k$ và khởi tạo.
  \item Thứ tự dữ liệu có thể ảnh hưởng kết quả cuối (với một số triển khai).
  \item \textbf{Nhạy với chuẩn hóa (rescale)}: chuẩn hóa/chuẩn tắc hóa làm thay đổi hoàn toàn phân cụm --- cần xử lý nhất quán.
  \item Kém khi cụm có \textbf{hình học phức tạp} (không lồi, không cầu).
\end{itemize}
\end{luuy}

\subsection{Ứng dụng K-Means}

\begin{ynghia}[title={Ứng dụng K-Means}]
  \begin{itemize}
    \item Phân đoạn ảnh: Gom vùng pixel theo màu hoặc texture; dùng trong nhận dạng đối tượng, truy vấn ảnh, ảnh y khoa.
    \item Phân khúc khách hàng: Chia khách theo hành vi mua hoặc nhân khẩu; hỗ trợ giữ chân khách, loyalty, quảng cáo nhắm mục tiêu.
    \item Phát hiện bất thường: Điểm không thuộc cụm nào (hoặc xa mọi centroid) có thể là bất thường (anomaly) --- gian lận, xâm nhập mạng, bảo trì dự đoán.
    \item Phân tích biểu hiện gen: Gom gen đồng điều hòa/biểu hiện; hỗ trợ khám phá thuốc, chẩn đoán, y học cá thể hóa.
  \end{itemize}
\end{ynghia}
